\section{Đánh giá thiết kế}

\subsection{Mục tiêu đánh giá}
Quá trình kiểm thử và đánh giá tính khả dụng thực nghiệm (Usability Testing) được tổ chức nhằm đo lường mức độ đáp ứng của bản mẫu đối với các mục tiêu trải nghiệm đã đặt ra \cite{nielsen1994usability, brooke1996sus}. Đợt đánh giá tập trung vào 5 khía cạnh then chốt:
\begin{enumerate}
    \item \textbf{Tính hiệu quả (Effectiveness)}: Đo lường tỷ lệ hoàn thành tác vụ thành công (Task Completion Rate).
    \item \textbf{Tính hiệu suất (Efficiency)}: Đo lường thời gian trung bình cần thiết để hoàn thành từng tác vụ (Time-on-Task).
    \item \textbf{Mức độ lỗi (Errors)}: Ghi nhận số lượng và phân loại các lỗi thao tác nhầm của người dùng.
    \item \textbf{Khả năng học hỏi (Learnability)}: Đánh giá độ trực quan của giao diện đối với người dùng tiếp cận lần đầu.
    \item \textbf{Mức độ hài lòng (Satisfaction)}: Đo lường cảm nhận chủ quan và mức độ an tâm qua thang đo Likert 5 mức độ.
\end{enumerate}

\subsection{Phương pháp đánh giá}
Phương pháp đánh giá chủ đạo là \textbf{Kiểm thử tính khả dụng có người điều phối (Moderated Usability Testing)} kết hợp \textbf{Giao thức suy nghĩ thành tiếng (Think-aloud Protocol)} trên thiết bị di động thực tế. Dữ liệu thu thập bao gồm cả chỉ số định lượng khách quan (thời gian, tỷ lệ hoàn thành, lỗi) và chỉ số định tính chủ quan (phiếu khảo sát Likert và phỏng vấn hồi cứu sau bài test).

\subsection{Người tham gia đánh giá}
Đợt kiểm thử thực nghiệm được thực hiện với \textbf{5 người dùng mục tiêu} đại diện đầy đủ cho 2 phân khúc Persona của dự án (dữ liệu chi tiết tại Bảng \ref{tab:test_participants}):

\begin{table}[H]
\centering
\small
\begin{tabularx}{\textwidth}{|c|p{2.8cm}|p{2.2cm}|X|p{3.2cm}|}
\hline
\textbf{Mã} & \textbf{Họ và Tên} & \textbf{Nhóm Persona} & \textbf{Nghề nghiệp / Bối cảnh} & \textbf{Thú cưng nuôi} \\ \hline
\textbf{P1} & Nguyễn Hoàng Lan & Persona 1 (Chính) & Chuyên viên Dữ liệu (28 tuổi) & Chó Poodle (Bơ, da nhạy cảm) \\ \hline
\textbf{P2} & Lê Thị Hồng Nhung & Persona 1 (Chính) & Nhân viên Văn phòng (26 tuổi) & Chó Corgi (Bánh Bao, dị ứng xà phòng) \\ \hline
\textbf{P3} & Phạm Văn Tuấn & Persona 1 (Chính) & Kỹ sư Phần mềm (30 tuổi) & Chó Poodle lai (Milo, da dễ ngứa) \\ \hline
\textbf{P4} & Trần Minh Khoa & Persona 2 (Phụ) & Sinh viên Năm 3 (21 tuổi) & Mèo Anh lông ngắn (Miu, nhút nhát) \\ \hline
\textbf{P5} & Nguyễn Phương Linh & Persona 2 (Phụ) & Sinh viên Năm 4 (22 tuổi) & Mèo Ba Tư (Bông, sợ tiếng ồn) \\ \hline
\end{tabularx}
\caption{Danh sách người dùng tham gia kiểm thử tính khả dụng thực nghiệm}
\label{tab:test_participants}
\end{table}

\subsection{Tác vụ kiểm thử}
Người tham gia được yêu cầu thực hiện 4 tác vụ cốt lõi dẫn xuất trực tiếp từ các kịch bản sử dụng thực tế:
\begin{itemize}
    \item \textbf{Tác vụ 1 (T1) --- Đặt lịch dịch vụ}: Thao tác chọn thú cưng, chọn gói dịch vụ tắm spa và chọn một khung giờ trống trong ma trận thời gian thực (Ngưỡng chuẩn kỳ vọng: $\le 45\text{s}$, Tỷ lệ hoàn thành $\ge 90\%$).
    \item \textbf{Tác vụ 2 (T2) --- Dặn dò dị ứng/y tế}: Kiểm tra và đính kèm ghi chú tiền sử dị ứng xà phòng và dặn dò thao tác nhẹ nhàng vào phiếu tiếp nhận (Ngưỡng chuẩn kỳ vọng: $\le 60\text{s}$, Tỷ lệ hoàn thành $\ge 85\%$).
    \item \textbf{Tác vụ 3 (T3) --- Theo dõi tiến độ Live Tracking}: Quan sát màn hình theo dõi tiến trình 4 mốc và nhận diện mốc trạng thái hiện tại của thú cưng (Ngưỡng chuẩn kỳ vọng: $\le 20\text{s}$, Tỷ lệ hoàn thành $\ge 95\%$).
    \item \textbf{Tác vụ 4 (T4) --- Tra cứu lịch sử \& Hóa đơn}: Tra cứu lại tên loại sữa tắm đã dùng ở lượt chăm sóc trước và xem chi tiết hóa đơn điện tử (Ngưỡng chuẩn kỳ vọng: $\le 35\text{s}$, Tỷ lệ hoàn thành $\ge 90\%$).
\end{itemize}

\subsection{Chỉ số đo lường}
\begin{itemize}
    \item \textbf{Tỷ lệ hoàn thành tác vụ (Completion Rate - \%)}: Tỷ lệ phần trăm người tham gia hoàn thành đúng mục tiêu tác vụ.
    \item \textbf{Thời gian thực hiện tác vụ (Time on Task - giây)}: Thời gian trung bình ($\bar{T}$) tính từ lúc bắt đầu đọc kịch bản đến khi hoàn tất thao tác cuối cùng.
    \item \textbf{Tần suất lỗi thao tác (Error Count)}: Số lần nhấp nhầm nút hoặc đi sai luồng màn hình.
    \item \textbf{Điểm hài lòng Likert (1--5)}: Thang điểm 5 mức độ đo lường tính trực quan, độ an tâm và sự tiện lợi sau bài test.
\end{itemize}

\subsection{Quy trình đánh giá}
Quy trình kiểm thử kéo dài 35 phút cho mỗi người tham gia, bao gồm 4 giai đoạn chuẩn mực:
\begin{enumerate}
    \item \textbf{Giai đoạn 1 --- Giới thiệu \& Hướng dẫn (5 phút)}: Giải thích mục đích kiểm thử, tạo tâm lý thoải mái và hướng dẫn phương pháp Think-aloud.
    \item \textbf{Giai đoạn 2 --- Thực thi nhiệm vụ (20 phút)}: Người tham gia lần lượt thực hiện 4 tác vụ độc lập trên điện thoại; điều phối viên bấm giờ và ghi nhật ký thao tác.
    \item \textbf{Giai đoạn 3 --- Điền phiếu khảo sát Likert (5 phút)}: Người tham gia trả lời 5 câu hỏi đánh giá trải nghiệm.
    \item \textbf{Giai đoạn 4 --- Phỏng vấn hồi cứu (5 phút)}: Phỏng vấn sâu về các điểm ngập ngừng hoặc khó khăn quan sát được.
\end{enumerate}

\subsection{Kết quả đánh giá}
Toàn bộ số liệu thực nghiệm được tổng hợp và phân tích chi tiết tại Bảng \ref{tab:test_performance_results} và Bảng \ref{tab:test_likert_results}:

\begin{table}[H]
\centering
\footnotesize
\begin{tabularx}{\textwidth}{|c|X|c|c|c|c|c|}
\hline
\textbf{Mã} & \textbf{Tên nhiệm vụ} & \textbf{Tỷ lệ thành công} & \textbf{Thời gian TB} & \textbf{Số lỗi TB} & \textbf{Ngưỡng chuẩn} & \textbf{Kết luận} \\ \hline
\textbf{T1} & Đặt lịch dịch vụ nhanh & \textbf{100\%} (5/5) & \textbf{39.0s} & 0.0 & $\le 45\text{s}$ ($\ge 90\%$) & \textbf{ĐẠT (PASS)} \\ \hline
\textbf{T2} & Dặn dò dị ứng/y tế & \textbf{100\%} (5/5) & \textbf{57.4s} & 1.0 & $\le 60\text{s}$ ($\ge 85\%$) & \textbf{ĐẠT (PASS)} \\ \hline
\textbf{T3} & Live Tracking 4 mốc & \textbf{100\%} (5/5) & \textbf{18.8s} & 0.2 & $\le 20\text{s}$ ($\ge 95\%$) & \textbf{XUẤT SẮC} \\ \hline
\textbf{T4} & Lịch sử \& Hóa đơn & \textbf{100\%} (5/5) & \textbf{30.8s} & 0.0 & $\le 35\text{s}$ ($\ge 90\%$) & \textbf{ĐẠT (PASS)} \\ \hline
\textbf{Tổng} & \textbf{Toàn hệ thống} & \textbf{100\%} (20/20) & --- & \textbf{0.30 lỗi/lượt} & $\ge 85\%$ & \textbf{XUẤT SẮC} \\ \hline
\end{tabularx}
\caption{Tổng hợp chỉ số hiệu năng thực nghiệm từ đợt kiểm thử khả dụng}
\label{tab:test_performance_results}
\end{table}

\begin{table}[H]
\centering
\small
\begin{tabularx}{\textwidth}{|c|X|c|c|}
\hline
\textbf{STT} & \textbf{Nội dung nhận định khảo sát} & \textbf{Điểm TB ($\bar{X}$)} & \textbf{Đánh giá} \\ \hline
\textbf{Q1} & Quy trình đặt lịch rõ ràng, trực quan và dễ thao tác & \textbf{4.40 / 5.0} & Rất cao \\ \hline
\textbf{Q2} & Thẻ cảnh báo dị ứng/dặn dò y tế nổi bật và tạo sự an tâm & \textbf{3.60 / 5.0} & Khá \\ \hline
\textbf{Q3} & Thanh tiến trình 4 mốc thời gian thực minh bạch, giảm lo âu & \textbf{4.20 / 5.0} & Rất cao \\ \hline
\textbf{Q4} & Tra cứu lịch sử dịch vụ và chi tiết sản phẩm thuận tiện & \textbf{4.40 / 5.0} & Rất cao \\ \hline
\textbf{Q5} & Mức độ hài lòng tổng thể và sẵn sàng sử dụng lâu dài & \textbf{4.20 / 5.0} & Rất cao \\ \hline
\textbf{Tổng} & \textbf{Điểm hài lòng trung bình toàn hệ thống} & \textbf{4.16 / 5.0} & \textbf{VƯỢT CHUẨN ($\ge 4.0$)} \\ \hline
\end{tabularx}
\caption{Kết quả khảo sát cảm nhận người dùng theo thang đo Likert 5 mức độ}
\label{tab:test_likert_results}
\end{table}

\subsection{Thảo luận và phân tích}
Mặc dù 100\% tác vụ đều đạt ngưỡng chuẩn kỳ vọng, quá trình quan sát thực nghiệm và phân tích định tính đã chỉ ra 4 phát hiện khả dụng (Usability Findings) cần cải tiến trước khi chuyển giao lập trình:

\begin{enumerate}
    \item \textbf{Phát hiện 1 (F1 --- Mức độ nghiêm trọng: Major)}: Ở tác vụ T2, người dùng P2 và P5 mất nhiều thời gian tìm ô nhập ghi chú do mục dặn dò nằm tách biệt trong tab Hồ sơ thú cưng. \textit{Giải pháp}: Tự động trích xuất ghi chú dị ứng từ hồ sơ và gắn Thẻ cảnh báo Amber Alert (\texttt{\#F59E0B}) kèm nút chỉnh sửa nhanh ngay trên form đặt lịch.
    \item \textbf{Phát hiện 2 (F2 --- Mức độ nghiêm trọng: Minor)}: Người dùng P2 ngập ngừng phân biệt giữa mốc “Đang chăm sóc” và “Hoàn tất” do thiếu tín hiệu thị giác động. \textit{Giải pháp}: Bổ sung hiệu ứng vòng tròn phát sáng chuyển động nhịp (Active Pulsing Beacon) tại mốc đang diễn ra.
    \item \textbf{Phát hiện 3 (F3 --- Mức độ nghiêm trọng: Minor)}: Viền ô chọn giờ trống hơi nhạt trong điều kiện ánh sáng mạnh. \textit{Giải pháp}: Tăng độ đậm viền (\texttt{border: 2px solid \#0D9488}) và thêm nhãn tóm tắt thời lượng dịch vụ.
    \item \textbf{Phát hiện 4 (F4 --- Mức độ nghiêm trọng: Cosmetic)}: Người dùng mong muốn lọc lịch sử theo thú cưng và tải hóa đơn PDF. \textit{Giải pháp}: Bổ sung thanh trượt Filter Chips và nút xuất hóa đơn PDF.
\end{enumerate}
